% LuaLaTeX 用。コンパイル例: lualatex swebok_ch07_software_maintenance_ja_regenerated.tex
\documentclass[a4paper,11pt]{ltjsarticle}
\usepackage[margin=22mm]{geometry}
\usepackage{array,longtable,booktabs}
\usepackage[unicode]{hyperref}
\usepackage{xcolor}
\usepackage{fancyhdr}
\usepackage{enumitem}
\hypersetup{colorlinks=true,linkcolor=black,urlcolor=black,citecolor=black}
\setlength{\parindent}{0pt}
\setlength{\parskip}{0.65em}
\setlist[itemize]{leftmargin=2em,itemsep=0.2em,topsep=0.2em}
\pagestyle{fancy}
\fancyhf{}
\lhead{07章 ソフトウェア保守}
\rhead{SWEBOK v4.0a 日本語まとめ}
\cfoot{\thepage}
\newcommand{\boxblock}[2]{\vspace{0.4em}\noindent\fbox{\begin{minipage}{0.96\linewidth}\textbf{#1}\par #2\end{minipage}}\vspace{0.4em}}
\newcommand{\todobox}[2]{\vspace{0.4em}\noindent\fbox{\begin{minipage}{0.96\linewidth}\textbf{TODO（図） #1}\par\textbf{画像生成用プロンプト}\par #2\end{minipage}}\vspace{0.4em}}
\begin{document}
\begin{titlepage}
\vspace*{20mm}
{\Huge\bfseries 07 章 ソフトウェア保守\par}
\vspace{6mm}
{\Large Software Maintenance の日本語まとめ\par}
\vspace{6mm}
SWEBOK Guide v4.0a Chapter 07 を初学者向けに再構成\par
\vfill
作成日：2026年5月12日\par
\end{titlepage}
\boxblock{この資料の主張}{ソフトウェア保守は、単なる「バグ修正」ではない。運用中のソフトウェアを価値ある状態に保ち、環境変化、利用者要求、障害、セキュリティ脅威、技術的陳腐化に対応し続ける活動である。保守は納品後だけに始まるものではなく、開発初期から保守性、移行、運用、支援、測定を計画する必要がある。}
\boxblock{この資料の読み方}{専門用語は原則として日本語で示し、必要に応じて英語を括弧内に併記する。英語名は、元資料やツール名を確認する手がかりとして残す。本文は初学者が前提知識なしに読めるように、概念の目的、実務で見るべき対象、用語の境界を重視している。}
この資料は SWEBOK Guide v4.0a Chapter 07 の内容を学習用に要約・再構成したものであり、逐語訳ではない。\par
\tableofcontents
\clearpage
\section{全体概要：この章で学ぶこと}
第7章は、運用中のソフトウェアを支え、変更し、価値を保つための知識領域である。ソフトウェアは納品された時点で完成して終わるわけではない。利用者の業務、法規制、外部システム、ハードウェア、クラウド基盤、セキュリティ状況は変化し続ける。したがって、ソフトウェアも変更され続ける必要がある。\par
\boxblock{到達目標}{この章を読み終えると、ソフトウェア保守の定義、保守が必要になる理由、保守の分類、技術的・管理的課題、保守プロセス、代表的な保守技法、保守ツールの役割を説明できるようになる。}
\begin{itemize}
\item ソフトウェア保守を、欠陥修正だけでなく、変更・適応・改善・支援・廃止まで含む活動として理解する。
\item 是正保守、予防保守、適応保守、追加保守、完全化保守、緊急保守を区別する。
\item プログラム理解、回帰テスト、影響分析、保守性、技術的負債の関係を説明する。
\item 保守の費用、測定、計画、構成管理、品質活動を実務上の管理対象として理解する。
\item 継続的統合、継続的デリバリ、継続的テスト、継続的デプロイが保守とどう関係するかを説明する。
\end{itemize}
\todobox{07 章の全体地図}{白背景、モノクロ、A4本文幅に収まる横長の学習用図解を作成する。中央に「ソフトウェア保守」を置き、周囲に「保守の基礎」「主要課題」「保守プロセス」「保守技法」「保守ツール」を配置する。各領域の下に代表語として「保守分類」「影響分析」「保守計画」「再工学」「静的解析」を小さく記載する。日本語ラベル、細線、講義ノート風、余白を広めにする。}
\subsection{最初に必要な用語}
\begin{longtable}{p{0.30\linewidth}p{0.64\linewidth}}
\toprule
用語 & 初学者向けの説明 \\
\midrule
\endhead
ソフトウェア保守（software maintenance） & 運用中のソフトウェアに対し、費用対効果の高い支援を提供するために必要な活動の総体である。修正、変更、監視、訓練、ヘルプデスク連携、廃止準備を含む。 \\
保守担当者（maintainer） & 保守活動を担う個人、チーム、または組織である。既存の成果物を理解し、変更し、運用中の品質を守る役割を持つ。 \\
変更要求（MR: Modification Request） & 機能追加、改善、設定変更など、ソフトウェア変更を求める依頼である。 \\
問題報告（PR: Problem Report） & 障害、不具合、利用上の問題、期待と実際の差異を報告する依頼である。 \\
保守性（maintainability） & ソフトウェアを修正、改善、適応、テストしやすい性質である。 \\
技術的負債（technical debt） & 短期的な都合で残された設計・実装・文書・テスト上の問題が、後の変更費用を増やす状態である。 \\
構成管理（SCM: Software Configuration Management） & ソフトウェア、文書、設定、テスト、リリースを識別し、変更を管理し、状態を記録する活動である。 \\
回帰テスト（regression testing） & 変更後に、以前は動いていた機能が壊れていないかを選択的に再テストすることである。 \\
サービス水準合意（SLA） & 保守や運用で提供するサービス内容と品質水準を、契約または合意として定めたものである。 \\
継続的統合（CI） & チームの変更を頻繁に共有の主線へ統合し、自動ビルドやテストで問題を早く見つける実践である。 \\
\bottomrule
\end{longtable}
\section{導入：ソフトウェア保守が重要な理由}
ソフトウェア開発が成功すると、利用者要求を満たすソフトウェア製品が提供される。しかし、運用が始まると、潜在していた欠陥が見つかり、運用環境が変化し、新しい利用者要求が現れる。したがって、ソフトウェア製品は進化しなければならない。\par
保守段階は、保証期間または導入後支援の終了後に始まると考えられがちである。しかし、実際の保守活動はもっと早く始まる。納品前から、運用後の支援、保守性、移行手順、支援体制を計画する必要がある。\par
近年は、開発、保守、運用の分断を減らし、継続的にソフトウェアを進化させる考え方が強くなっている。DevOps の実践やツールが普及したことで、開発者、運用担当者、保守担当者が協力し、素早く安全に変更を届けることが重視されている。\par
\boxblock{重要}{保守は「古いソフトウェアを仕方なく直す作業」ではない。運用中のソフトウェアの価値を維持し、変化に対応し、利用者に継続的な便益を届けるための中核的なソフトウェア工学活動である。}
\section{1 ソフトウェア保守の基礎}
この節では、ソフトウェア保守の役割と範囲を理解するための概念と用語を整理する。特に重要なのは、保守を種類別に分類して理解することである。分類ができると、費用、優先順位、必要な人員、テスト範囲、リリース判断を説明しやすくなる。\par
\subsection{1.1 定義と用語}
ソフトウェア保守の目的は、既存ソフトウェアを変更しながら、その完全性を保つことである。完全性とは、単に動くことではなく、要求、設計、コード、テスト、文書、構成、運用上の整合が保たれている状態を意味する。\par
SWEBOK では、ソフトウェア保守を「運用中のソフトウェアに費用対効果の高い支援を提供するために必要な活動の総体」として扱う。したがって、保守には納品後の修正だけでなく、納品前の保守計画、保守性の確保、移行支援、訓練、ヘルプデスク連携も含まれる。\par
\boxblock{定義}{ソフトウェア保守とは、運用中のソフトウェアを支援し、変更し、品質と価値を保つための活動である。変更しながら壊さないこと、すなわち「変化」と「完全性」の両立が中心課題である。}
\subsection{1.2 ソフトウェア保守の性質}
保守は、開発から運用、最終的な廃止まで、ソフトウェアのライフサイクル全体を支える。運用中のソフトウェアは、容量、継続性、可用性の観点で監視される。変更要求や問題報告は記録され、追跡され、影響が分析される。その後、コードや文書などの成果物が修正され、テストされ、新しい版が運用へリリースされる。\par
保守担当者は、開発者や運用者が持つ知識から学ぶ必要がある。保守担当者が早期に開発へ関与すれば、保守費用と手戻りを減らしやすい。反対に、開発者が離れた後に保守担当者が参加すると、コード、テスト、文書、設計意図を理解する負担が大きくなる。\par
\subsection{1.3 ソフトウェア保守が必要な理由}
ソフトウェア保守は、ソフトウェアが寿命の間ずっと利用者要求を満たし続けるために必要である。開発ライフサイクルがウォーターフォール型でもアジャイル型でも、ソフトウェアは変更される。変更の理由は、欠陥修正だけではない。\par
\begin{itemize}
\item 欠陥や潜在欠陥を修正するため。
\item 運用中ソフトウェアの設計、性能、信頼性を改善するため。
\item 新しい機能や利用者要求を実装するため。
\item 利用者が機能を理解し、使い続けられるように支援するため。
\item 外部システム、基盤、ハードウェア、クラウド、法規制の変化へ適応するため。
\item セキュリティ脅威を予防し、脆弱性を修正するため。
\item 技術的陳腐化に対応し、老朽化した要素を更新するため。
\item 不要になったソフトウェアを安全に廃止するため。
\end{itemize}
\subsection{1.4 保守費用が大きい理由}
ソフトウェアの総費用の中で、保守は大きな割合を占める。一般に、誤りの修正費用はライフサイクルの後半ほど大きくなる。また、保守は長期間続くため、単発の開発費用よりも見えにくいが、累積費用は大きい。\par
保守は「欠陥修正が中心」と誤解されやすい。しかし、調査では保守の多くが機能強化や環境適応に使われるとされる。修正と強化を同じ報告カテゴリにまとめると、欠陥修正が実際以上に高く見える。保守分類を分けて測定することは、費用構造を理解するために重要である。\par
\boxblock{費用を見る観点}{保守費用は、運用環境と組織環境の両方に左右される。運用環境にはハードウェア、OS、ミドルウェア、クラウド基盤が含まれる。組織環境には方針、競争、プロセス、製品、人員が含まれる。}
\subsection{1.5 ソフトウェアの進化}
ソフトウェア保守は、ソフトウェア進化を支える活動である。Lehman らの研究は、長期運用されるソフトウェアが変化し続けること、そして変化を管理しなければ複雑さや品質低下が進むことを示した。\par
\begin{longtable}{p{0.30\linewidth}p{0.64\linewidth}}
\toprule
進化の法則 & 初学者向けの意味 \\
\midrule
\endhead
継続的変化 & ソフトウェアは利用環境に合わせて継続的に適応しなければ、満足度が下がる。 \\
複雑さの増大 & 変更を重ねると、複雑さは増える。複雑さを減らす作業をしなければ保守は難しくなる。 \\
自己調整 & 進化の過程は、製品やプロセスの測定値を通じて一定の傾向を持つ。 \\
作業率の不変性 & 長期的には、進化するソフトウェアに投入できる有効作業量には一定の制約がある。 \\
親熟度の保存 & 関係者がソフトウェアの内容と振る舞いを理解し続ける必要がある。成長が急すぎると理解が追いつかない。 \\
継続的成長 & 利用者満足を保つため、機能内容は継続的に増える傾向がある。 \\
品質低下 & 運用環境の変化へ厳密に適応しなければ、品質は低下して見える。 \\
フィードバックシステム & ソフトウェア進化は複数の利害関係者と情報の循環から成るため、フィードバックとして扱う必要がある。 \\
\bottomrule
\end{longtable}
\todobox{ソフトウェア進化の循環図}{白背景、モノクロ、A4本文幅に収まる循環図を作成する。「環境変化」「要求追加」「ソフトウェア変更」「複雑さ増大」「保守性低下」「再整理・リファクタリング」を円環状に配置する。中央に「ソフトウェア進化」と書く。矢印は細線、講義ノート風、日本語ラベル。}
\subsection{1.6 ソフトウェア保守の分類}
ソフトウェア保守は、変更要求を分類するために複数の種類へ分けられる。SWEBOK では、是正、予防、適応、追加、完全化に加えて、緊急保守を説明する。\par
\begin{longtable}{p{0.22\linewidth}p{0.40\linewidth}p{0.30\linewidth}}
\toprule
分類 & 説明 & 例 \\
\midrule
\endhead
是正保守 & 納品後に見つかった問題を修正する反応的な変更である。 & 本番障害の原因となった計算誤りを修正する。 \\
予防保守 & 本番で問題が起きる前に潜在欠陥を修正する変更である。 & 将来エラーになり得る古い日付処理を直す。 \\
適応保守 & 環境変化の中でソフトウェアを使い続けられるようにする変更である。 & OS、データベース、外部APIの変更に対応する。 \\
追加保守 & 利用を拡張するために、比較的大きな機能や特徴を追加する変更である。 & 既存サービスに新しい通知機能を追加する。 \\
完全化保守 & 利用者向け改善、文書改善、性能や保守性などの品質改善を行う変更である。 & 画面操作を簡単にする、コードを読みやすくする。 \\
緊急保守 & 是正保守までの間、一時的にシステムを稼働させ続けるための予定外変更である。 & 本番停止を避けるため、一時的な回避策を適用する。 \\
\bottomrule
\end{longtable}
\boxblock{分類の使い方}{保守分類は、会計上の分類ではなく、意思決定の道具である。どの保守が多いかを測れば、品質問題、環境変化、利用者要求、技術的負債のどこに費用が使われているかが見える。}
\todobox{ソフトウェア保守分類図}{白背景、モノクロ、A4本文幅に収まる階層図を作成する。最上位に「変更要求」。次の層に「修正」と「強化」を置く。「修正」から「是正保守」「予防保守」へ分岐し、「強化」から「適応保守」「追加保守」「完全化保守」へ分岐する。横に「緊急保守」を点線枠で置き、一時対応であることを注記する。日本語ラベル、細線、講義ノート風。}
\section{2 ソフトウェア保守の主要課題}
保守には、既存ソフトウェアを変更するという独自の難しさがある。保守担当者は、しばしば他人が開発した大規模で複雑なソフトウェアの中から原因を見つけ、変更し、予期しない影響を避ける必要がある。\par
\begin{itemize}
\item 技術的課題：限られた理解、テスト、影響分析、保守性。
\item 管理上の課題：組織目標との整合、人員配置、プロセス、供給者管理、組織構造。
\item 費用の課題：技術的負債と保守費用の見積り。
\item 測定の課題：どこに労力と費用が使われているかを見える化すること。
\end{itemize}
\subsection{2.1 技術的課題}
\subsubsection{2.1.1 限られた理解}
限られた理解とは、保守担当者が、他者によって開発されたソフトウェアを最初は十分に理解していない状態である。変更すべき場所を見つけ、何を直せばよいかを理解するまでに大きな時間がかかる。\par
理解を難しくする要因には、領域知識の不足、読みにくい実装、文書不足、構成や提示の悪さがある。ソースコードのようなテキスト中心の表現では、変更履歴や設計意図を追うことが難しい。コード可視化、リバースエンジニアリング、明確な文書、コードブラウザが理解を助ける。\par
\subsubsection{2.1.2 テスト}
保守では、変更要求や問題報告の処理中にテスト計画とテスト実行が行われる。大規模ソフトウェアで完全な再テストを繰り返す費用は高い。そのため、変更に応じて適切なテストを計画し、特に回帰テストを重視する必要がある。\par
回帰テストは、変更が意図しない副作用を起こしていないかを確認する選択的な再テストである。保守では複数の担当者が同時に異なる問題へ対応するため、テスト調整も難しくなる。さらに、重要機能を持つ本番システムを停止してテストすることが難しい場合もある。\par
\subsubsection{2.1.3 影響分析}
影響分析とは、提案された変更が既存ソフトウェアへ与える詳細な影響を評価する活動である。変更対象、影響を受ける構成要素、必要な資源、リスク、必要なテスト、性能・安全性・セキュリティへの影響を調べる。\par
影響分析では、変更要求や問題報告をまずソフトウェアの言葉へ翻訳する必要がある。たとえば「画面が遅い」という報告は、データベース照会、ネットワーク、キャッシュ、画面描画、外部サービスなど、複数の技術要素へ分解して考える必要がある。\par
\begin{longtable}{p{0.30\linewidth}p{0.64\linewidth}}
\toprule
影響分析で決めること & 説明 \\
\midrule
\endhead
識別方法 & 変更要求や問題報告を一意に識別する方法を決める。 \\
分類と優先度 & 緊急度、重要度、保守分類、リスクで並べる。 \\
提出手順 & 利用者や運用者がどのように要求・報告を出すかを決める。 \\
追跡情報 & 利用者へ何を報告し、どの測定値で状況を示すかを決める。 \\
暫定回避策 & 本修正までの間、どのような回避策を提供するかを決める。 \\
フィードバック & 利用者へどの段階で何を返すかを決める。 \\
\bottomrule
\end{longtable}
\todobox{影響分析の連鎖図}{白背景、モノクロ、A4本文幅の横長図を作成する。左から「変更要求・問題報告」「要求」「設計」「コード」「テスト」「運用手順」「利用者文書」を箱で並べ、影響が伝播する矢印を描く。変更対象から影響範囲を点線で囲む。日本語ラベル、講義ノート風。}
\subsubsection{2.1.4 保守性}
保守性とは、ソフトウェア製品を修正できる能力である。修正には、欠陥の修正、改善、環境変化への適応、要求や機能仕様の変更が含まれる。保守性は重要な品質特性であり、開発中に仕様化、レビュー、管理されるべきである。\par
保守性が低いと、将来の保守担当者の負担が増え、技術的負債が蓄積する。緊急対応や短期的な機能追加で、十分なレビューや設計配慮を省くと、後からより多くの時間と労力が必要になる。\par
\boxblock{技術的負債を見る三つの観点}{第一に、コード品質と業務上の重要度を分ける。すべての負債が同じ緊急度ではない。第二に、アーキテクチャが組織目標に合っているかを見る。第三に、チームやプロセス上の欠落が負債を増やしていないかを見る。}
\subsection{2.2 管理上の課題}
\subsubsection{2.2.1 組織目標との整合}
保守活動は、事業、顧客、利用者の優先順位と整合していなければならない。新規開発は、期限と予算を持つプロジェクトとして見えやすい。一方、保守は長期的に運用ソフトウェアの寿命を延ばし、価値を保つ活動であるため、経済効果が見えにくい。\par
その結果、リファクタリング、セキュリティ改善、性能改善よりも、目に見える新機能が優先されることがある。しかし、新機能だけを優先すると、コード劣化と技術的負債が増え、長期的な変更速度が落ちる。組織目標、保守性、工学標準の均衡が必要である。\par
\begin{itemize}
\item 運用中ソフトウェアの全ライフサイクル費用を理解する。
\item 新規開発と既存ソフトウェア改善の費用対効果を比較する。
\item 同じチームが開発と保守に責任を持つことで、技能と責任の分断を減らす。
\item 初期開発から保守性要求に注意を向ける。
\end{itemize}
\subsubsection{2.2.2 人員配置}
保守作業は、しばしば新規開発より魅力が低いと見られる。しかし、保守はライフサイクル活動の大きな部分を占め、保守担当者の貢献は重要である。組織は、保守担当者の専門性を評価し、学習機会とキャリア上の意味を与える必要がある。\par
人員配置では、既存コードを読み解く能力、運用への理解、テスト設計、障害対応、コミュニケーション能力が必要である。保守担当者が孤立すると、知識が属人化し、対応品質が下がる。\par
\subsubsection{2.2.3 プロセス}
ソフトウェアライフサイクルプロセスは、ソフトウェアと関連成果物を開発・保守するための活動、方法、実践、変換の集合である。保守は開発と多くを共有する。たとえば、構成管理、変更管理、レビュー、テスト、文書更新はどちらにも必要である。\par
一方で、保守には開発にない活動もある。既存プログラム理解、運用からの問題報告処理、ヘルプデスク連携、緊急修正、リリース後の小さな変更群の管理がその例である。\par
\subsubsection{2.2.4 供給者管理}
供給者管理とは、保守を外部供給者に委託する場合に、供給者とその成果を適切に管理することである。外部委託、オフショア保守、複数供給者によるサービス提供は、現代の保守では一般的な選択肢である。\par
\begin{longtable}{p{0.30\linewidth}p{0.64\linewidth}}
\toprule
契約形態 & 説明 \\
\midrule
\endhead
単一供給者・提携 & 一つの供給者がすべての保守サービスを提供する。または外部の統合者が複数供給者との関係を管理する。 \\
複数供給者 & 複数の独立した供給者が製品やサービスを提供し、それらを単一のソフトウェア支援サービスとして統合する。XaaS の普及により増えている。 \\
\bottomrule
\end{longtable}
外部委託では、保守範囲、サービス水準合意、契約上の責任、連絡体制、ヘルプデスク、プロセス自動化を明確にする必要がある。特に中核業務を支えるソフトウェアでは、組織が制御を失わないように注意する。\par
\subsubsection{2.2.5 保守の組織面}
保守組織の設計では、誰が保守に責任を持つかを決める。アジャイル型では、開発者が保守も担い、開発者兼保守担当者として働くことが多い。一方で、開発チームとは別の保守チームを置く組織もある。\par
\begin{longtable}{p{0.22\linewidth}p{0.40\linewidth}p{0.30\linewidth}}
\toprule
方式 & 利点 & 注意点 \\
\midrule
\endhead
同一チーム方式 & 開発知識が残り、要求変化へ素早く対応しやすい。 & 障害対応が新規開発を中断しやすい。文書化が軽視されることがある。 \\
分離チーム方式 & 運用と保守に集中でき、長期支援を組織化しやすい。 & 引き継ぎ摩擦、知識喪失、保守担当者の動機低下が起こり得る。 \\
\bottomrule
\end{longtable}
\boxblock{判断基準}{どの組織形態が正しいかは一律には決まらない。重要なのは、経験ある人やチームへ保守責任を委ね、保守作業と変更履歴に関する品質の高い文書を残すことである。}
\subsection{2.3 保守費用}
保守費用を理解するには、保守分類ごとに費用傾向を見る必要がある。是正、予防、適応、追加、完全化のどこに努力が使われているかを示せば、顧客や組織は保守の価値と課題を理解しやすくなる。\par
\subsubsection{2.3.1 技術的負債の費用見積り}
技術的負債は、コードを必要以上に保守しにくくし、変更費用を増やす。技術的負債の費用を見積もるには、規模、複雑さ、設計原則違反、コーディング慣行違反、コードのにおい、文書不足、テスト不足などを見る。\par
保守性を改善するには、読みやすさを確保し、構造化されたコードを保ち、複雑さを減らし、正確なコメントを付け、言語上の弱点やコンパイラ依存構造を避け、エラー追跡を容易にし、設計とコードの追跡可能性を保ち、コードレビューを行う必要がある。\par
\todobox{技術的負債と保守費用の関係図}{白背景、モノクロ、A4本文幅の概念図を作成する。横軸に「時間」、縦軸に「変更費用」。技術的負債を放置した曲線は右上がりに急増、継続的な改善を行った曲線は緩やかに増えるように描く。途中に「短期対応」「複雑化」「理解コスト増」「リファクタリング」を注記する。日本語ラベル、講義ノート風。}
\subsubsection{2.3.2 保守費用見積り}
保守費用の見積りは、ソフトウェア計画の早い段階で作成すべきである。見積りは保守活動の範囲に依存する。保守が進むにつれて、過去の測定データを使って見積りを更新する。\par
\begin{longtable}{p{0.30\linewidth}p{0.64\linewidth}}
\toprule
費用要素 & 説明 \\
\midrule
\endhead
利用者拠点への移動 & 現地支援や調査が必要な場合の費用である。 \\
訓練 & 保守担当者と利用者への教育費用である。 \\
開発・テスト環境 & ソフトウェア工学環境、テスト環境、ツールの維持費である。 \\
人件費 & 給与、福利厚生、要員確保にかかる費用である。 \\
消耗品・設備 & 保守作業に必要な追加資源である。 \\
ライセンス保守 & 商用ソフトウェアや基盤ライセンスの維持費である。 \\
製品変更と管理 & リリース管理、プログラム管理、変更実施にかかる費用である。 \\
\bottomrule
\end{longtable}
個別の変更要求や問題報告についても、影響分析の中で費用見積りが必要である。見積りには、実施工数、必要資源、実施期間、リスクを含める。パラメトリックモデル、類似システム比較、経験データなどの方法を使う。\par
\subsection{2.4 ソフトウェア保守の測定}
測定可能な保守対象には、保守プロセス、資源、製品がある。測定値には、規模、複雑さ、品質、理解しやすさ、保守性、労力が含まれる。特に、保守分類ごとに使われた労力を測ると、費用構造が見える。\par
\begin{longtable}{p{0.30\linewidth}p{0.64\linewidth}}
\toprule
測定対象 & 説明 \\
\midrule
\endhead
保守性 & モジュール性、再利用性、解析容易性、修正容易性、テスト容易性、支援容易性などで測る。 \\
信頼性 & 成熟性、可用性、耐故障性、回復性などで測る。 \\
規模 & 機能規模、コード行数、構成要素数などで測る。 \\
要求数 & 期間ごとの変更要求数、問題報告数、分類別件数を見る。 \\
労力 & 一件あたりの保守工数、分類別工数、リリースあたり工数を見る。 \\
技術特性 & プラットフォーム、言語、フレームワーク、ハードウェアなどを記録する。 \\
\bottomrule
\end{longtable}
\boxblock{測定の目的}{測定は管理のためだけではない。保守費用がどこに使われているかを説明し、改善対象を見つけ、組織内でアプリケーションごとの保守状況を比較するために使う。}
\section{3 ソフトウェア保守プロセス}
ソフトウェア保守には、ISO/IEC/IEEE 14764 などに示される標準的な活動に加え、保守担当者に固有の活動がある。保守は、ISO/IEC/IEEE 12207 で扱われる技術的ライフサイクルプロセスの一つでもある。\par
\subsection{3.1 ソフトウェア保守プロセス}
保守プロセスは、運用中ソフトウェアを支援するための活動、入力、出力を提供する。代表的なプロセスには、保守準備、保守実施、物流的支援、保守と支援結果の管理がある。\par
\begin{longtable}{p{0.30\linewidth}p{0.64\linewidth}}
\toprule
プロセス & 説明 \\
\midrule
\endhead
保守準備 & 保守体制、手順、環境、契約、知識移転を整える。 \\
保守実施 & 変更要求や問題報告を処理し、修正、テスト、リリースを行う。 \\
物流的支援 & ソフトウェア配布、導入、資源、支援物品、環境を扱う。 \\
結果管理 & 保守結果、支援状況、顧客満足、測定値、改善点を管理する。 \\
\bottomrule
\end{longtable}
近年は、軽量なアジャイル手法も保守に適用される。利用者が保守サービスに速い応答を求めるためである。保守プロセスの改善には、保守成熟度モデルや継続的改善の考え方が役立つ。\par
\todobox{ソフトウェア保守プロセス図}{白背景、モノクロ、A4本文幅のプロセス図を作成する。外側に「開発」「運用」「移行」「廃止」を配置し、中央に「保守準備」「保守実施」「物流的支援」「保守・物流結果の管理」を箱で置く。各箱を矢印でつなぎ、運用中ソフトウェアを支える流れを示す。日本語ラベル、細線、講義ノート風。}
\subsection{3.2 ソフトウェア保守活動とタスク}
保守プロセスには、既存ソフトウェアシステムを完全性を保ちながら運用・変更するための活動とタスクが含まれる。保守担当者は、開発と同様に、分析、設計、コーディング、テスト、文書化を行う。要求を追跡し、基準となる成果物が変われば文書も更新する。\par
\begin{longtable}{p{0.30\linewidth}p{0.64\linewidth}}
\toprule
保守に特徴的な活動 & 説明 \\
\midrule
\endhead
プログラム理解 & ソフトウェアが何を行い、各部分がどう連携するかを理解する活動である。 \\
移行 & 開発チームから運用・保守チームへ、ソフトウェアを段階的に移す管理された活動である。 \\
変更要求の受理・却下 & 合意された規模、労力、複雑さを超える依頼を、保守ではなく開発へ戻す判断を含む。 \\
保守ヘルプデスク & 利用者支援と保守調整を行い、要求や障害の評価、優先度付け、費用見積りを起動する。 \\
影響分析 & 変更が影響する領域を特定する活動である。 \\
サービス水準管理 & SLI、SLO、SLA、ライセンス、契約を通じて、支援内容と品質目標を管理する。 \\
\bottomrule
\end{longtable}
\subsubsection{3.2.1 支援・監視活動}
保守担当者は、文書化、構成管理、検証と妥当性確認、問題解決、ソフトウェア品質保証、レビュー、脆弱性評価、監査などの継続的支援活動を行う。顧客満足の監視も、保守結果の管理として重要である。\par
\subsubsection{3.2.2 計画活動}
保守では計画が重要である。開発プロジェクトは数か月から数年で終わることが多いが、保守はソフトウェアが廃止されるまで続く。したがって、保守計画は新しいソフトウェアを開発する判断の時点から始めるべきである。\par
\begin{longtable}{p{0.30\linewidth}p{0.64\linewidth}}
\toprule
計画の階層 & 説明 \\
\midrule
\endhead
事業計画 & 組織レベルで、予算、顧客、サービス方針、保守体制を決める。 \\
保守計画 & 移行レベルで、保守範囲、プロセス、ツール、体制、費用を決める。 \\
リリース・版計画 & ソフトウェアレベルで、次のリリース内容、日程、リスク、撤退計画を決める。 \\
変更要求計画 & 個別要求レベルで、影響分析、実施工数、優先度、完了時期を決める。 \\
\bottomrule
\end{longtable}
\todobox{保守計画の四層図}{白背景、モノクロ、A4本文幅のピラミッド図または層構造図を作成する。下から「個別変更要求計画」「リリース・版計画」「保守計画」「事業計画」を積み上げる。右側に「粒度は下ほど細かい」「期間は上ほど長い」と注記する。日本語ラベル、講義ノート風。}
\subsubsection{3.2.3 構成管理}
構成管理は、保守プロセスを支えるシステムまたはサービスである。変更の識別、承認、実装、リリース、検証、妥当性確認、監査を管理する。保守では、変更要求や問題報告だけでなく、ソフトウェア製品、その基盤、文書、テスト、設定も管理しなければならない。\par
保守における構成管理は、運用環境で発生する多数の小さな変更を統制する点で難しい。保守担当者は構成管理計画と運用手順に従い、構成管理委員会などで次のリリース内容を判断する。\par
\subsubsection{3.2.4 ソフトウェア品質}
保守を行えば自動的に品質が上がるわけではない。保守担当者は有効な品質プログラムを持つ必要がある。品質保証、検証と妥当性確認、レビュー、監査の活動と技法は、他のプロセスと整合させて選択する。\par
運用と保守は、開発プロセスの成果物を再利用すべきである。たとえば、テストツール、文書、テスト結果は保守にも役立つ。開発時に保守を意識して成果物を残すことが、長期的な品質を支える。\par
\section{4 ソフトウェア保守技法}
この節では、保守で広く用いられる技法を整理する。技法は、既存ソフトウェアを理解し、変更し、影響を抑え、継続的に品質を守るために使われる。\par
\subsection{4.1 プログラム理解}
プログラム理解とは、変更を実装するために、既存プログラムを読み、構造と振る舞いを理解する活動である。保守担当者は、コードの検索、参照、依存関係の確認、実行経路の確認を行う。コードブラウザや明確な文書は、理解を大きく助ける。\par
\boxblock{実務の目安}{保守作業では、いきなりコードを変えるのではなく、まず対象機能、関連モジュール、依存する外部サービス、テスト、設定、過去の変更履歴を確認する。}
\subsection{4.2 ソフトウェア再工学}
ソフトウェア再工学とは、既存ソフトウェアを調査し、変更し、新しい形に再構成することである。老朽化したレガシーソフトウェアを置き換える目的で行われることが多い。\par
リファクタリングは、外部から見た振る舞いを変えずに、内部構造を整理する再工学技法である。保守活動の前後にリファクタリングを行うと、技術的負債を減らし、将来の変更を容易にできる。\par
アジャイル開発や継続的統合では、小さな変更が頻繁に行われるため、コード品質と信頼性を保つために継続的なリファクタリングが重要である。\par
\subsection{4.3 リバースエンジニアリング}
リバースエンジニアリングとは、既存ソフトウェアを分析して、構成要素と相互関係を明らかにし、より抽象度の高い表現を作る活動である。リバースエンジニアリング自体は受動的であり、ソフトウェアを変更しない。\par
\begin{longtable}{p{0.30\linewidth}p{0.64\linewidth}}
\toprule
種類 & 説明 \\
\midrule
\endhead
再文書化 & 既存コードや設定から、理解しやすい文書や図を作り直す。 \\
設計回復 & コードから設計意図、構造、アーキテクチャを復元する。 \\
データ・リバースエンジニアリング & 物理データベースから論理スキーマやデータ関係を回復する。 \\
\bottomrule
\end{longtable}
リバースエンジニアリングでは、ツールが重要である。呼出しグラフ、制御フローグラフ、依存関係図、データフロー図などを生成し、保守担当者が構造と変化を理解する助けになる。\par
\subsection{4.4 継続的統合・継続的デリバリ・継続的テスト・継続的デプロイ}
開発、運用、保守に関する作業を自動化すると、工学的資源を節約できる。適切に実装された自動化は、手作業より速く、容易で、信頼しやすい。DevOps は、開発、運用、保守の資源と手順を組み合わせ、継続的統合、継続的デリバリ、継続的テスト、継続的デプロイを実行する。\par
\begin{longtable}{p{0.30\linewidth}p{0.64\linewidth}}
\toprule
実践 & 説明 \\
\midrule
\endhead
継続的統合（CI） & チーム全員の変更を頻繁に共有の主線へ統合し、自動ビルドとテストで統合誤りを早く検出する。 \\
継続的デリバリ & 最新のコードと設定を継続的に組み立て、ステージング環境やテスト環境へ頻繁にリリース候補を届ける。 \\
継続的テスト & ライフサイクルの各段階でテストを行い、早く頻繁に品質を評価する。 \\
継続的デプロイ & 検証済みの変更を本番環境へ自動的に展開し、リスクを抑えながら価値を届ける。 \\
\bottomrule
\end{longtable}
\todobox{CI/CD/テスト/デプロイのパイプライン図}{白背景、モノクロ、A4本文幅の横長パイプライン図を作成する。左から「変更」「継続的統合」「自動ビルド」「継続的テスト」「継続的デリバリ」「継続的デプロイ」「運用フィードバック」を並べ、矢印でつなぐ。上に「保守変更を早く安全に届ける」と注記する。日本語ラベル、講義ノート風。}
\subsection{4.5 保守の可視化}
保守では、進化し続ける構造と依存関係を理解し続けることが難しい。可視化は、ソフトウェアの構成要素、依存関係、変更履歴、実行時の振る舞いを視覚的に表し、保守担当者の判断を助ける。\par
大規模で複雑なソフトウェアでは、可視化により依存関係分析、進化履歴の追跡、実行時動作の理解、補完的な文書化が可能になる。可視化は、計算機の処理能力と人間のパターン認識能力を組み合わせる技法である。\par
\todobox{保守可視化ダッシュボード}{白背景、モノクロ、A4本文幅に収まるダッシュボード風図を作成する。画面内に「依存関係図」「変更履歴タイムライン」「技術的負債指標」「保守要求件数」「テスト状況」を小さなパネルとして配置する。日本語ラベル、細線、講義ノート風。}
\section{5 ソフトウェア保守ツール}
保守ツールは、既存ソフトウェアを変更する場面で特に重要である。保守ツールは開発ツールや運用ツールと密接に結びついており、全体としてソフトウェア工学環境の一部を構成する。\par
\begin{longtable}{p{0.30\linewidth}p{0.64\linewidth}}
\toprule
ツール分類 & 主な役割 \\
\midrule
\endhead
構成管理・版管理・コードレビュー & 変更を記録し、承認し、リリース内容を管理する。 \\
ソフトウェアテストツール & 単体テスト、結合テスト、回帰テスト、自動テストを支援する。 \\
品質評価ツール & 技術的負債、コード品質、複雑度、コードのにおいを評価する。 \\
プログラムスライサ & 変更の影響を受けるプログラム部分だけを選び出す。 \\
静的解析器 & プログラム内容を実行せずに確認し、構造、欠陥、脆弱性を調べる。 \\
動的解析器 & 実行時の経路や状態を追跡し、実際の振る舞いを確認する。 \\
データフロー解析器 & プログラム内のデータの流れを追跡する。 \\
相互参照生成器 & プログラム構成要素の索引を作る。 \\
依存関係解析器 & 構成要素同士の関係を理解する。 \\
遠隔アクセスツール & 利用者環境を遠隔で診断し、必要な修正を行う。 \\
リバースエンジニアリングツール & 既存製品から仕様、設計記述、図などの成果物を作る。 \\
\bottomrule
\end{longtable}
\todobox{保守ツールの分類図}{白背景、モノクロ、A4本文幅の分類図を作成する。中央に「保守ツール」を置き、周囲に「変更管理」「テスト」「品質評価」「静的解析」「動的解析」「依存関係解析」「リバースエンジニアリング」「遠隔支援」を放射状に配置する。各項目に短い説明を添える。日本語ラベル、講義ノート風。}
\section{6 トピックと参考文献の対応}
この表は、SWEBOK が示すトピック対参考資料の行列を、学習用に簡略化したものである。第7章では、Grubb and Takang の Software Maintenance: Concepts and Practice が中心的な参考資料として扱われる。\par
\begin{longtable}{p{0.30\linewidth}p{0.64\linewidth}}
\toprule
トピック & 主な参考資料 \\
\midrule
\endhead
1. ソフトウェア保守の基礎 & ISO/IEC/IEEE 14764; Grubb and Takang \\
1.1 定義と用語 & ISO/IEC/IEEE 14764; Grubb and Takang \\
1.2 保守の性質 & Grubb and Takang \\
1.3 保守が必要な理由 & Grubb and Takang \\
1.4 保守費用の大きさ & Grubb and Takang; Abran and Nguyenkim \\
1.5 ソフトウェア進化 & Lehman; Grubb and Takang \\
1.6 保守分類 & ISO/IEC/IEEE 14764; Grubb and Takang \\
2.1 技術的課題 & Grubb and Takang; ISO/IEC/IEEE 14764 \\
2.1.1 限られた理解 & Grubb and Takang \\
2.1.2 テスト & Grubb and Takang; Software Testing KA \\
2.1.3 影響分析 & ISO/IEC/IEEE 14764; Grubb and Takang \\
2.1.4 保守性 & ISO/IEC/IEEE 14764; Software Quality KA \\
2.2 管理上の課題 & Grubb and Takang; April and Abran \\
2.3 保守費用 & ISO/IEC/IEEE 14764; Grubb and Takang \\
2.4 保守測定 & Grubb and Takang; Software Engineering Management KA \\
3. 保守プロセス & ISO/IEC/IEEE 14764; ISO/IEC/IEEE 12207; Grubb and Takang \\
4. 保守技法 & Grubb and Takang; Humble and Fairley; Software Quality KA \\
5. 保守ツール & Grubb and Takang; Software Configuration Management KA \\
\bottomrule
\end{longtable}
\section{7 発展的な読み物}
\begin{longtable}{p{0.30\linewidth}p{0.64\linewidth}}
\toprule
文献 & 読むと得られること \\
\midrule
\endhead
April and Abran, Software Maintenance Management & 継続的な保守プロセス改善、成熟度、ベンチマーク、標準との整合を学べる。 \\
IEEE Std 2675-2021, DevOps & 信頼性とセキュリティを考慮したビルド、パッケージ、デプロイの原則とプロセスを学べる。 \\
Humble and Farley, Continuous Delivery & ビルド、テスト、デプロイの自動化を通じて、安全に頻繁なリリースを行う考え方を学べる。 \\
Lehman, Laws of Software Evolution & 長期運用されるソフトウェアがどのように変化し、複雑化するかを理解できる。 \\
Grubb and Takang, Software Maintenance & ソフトウェア保守の概念、実践、技法、測定を体系的に学べる。 \\
\bottomrule
\end{longtable}
\section{8 参考文献}
\begin{itemize}
\item IEEE Std, ISO/IEC/IEEE 14764:2022, Software Engineering — Software Life Cycle Processes — Maintenance, third edition.
\item P. Grubb and A. A. Takang, Software Maintenance: Concepts and Practice, 2nd ed., World Scientific Publishing, 2003.
\item A. April and A. Abran, Software Maintenance Management: Evaluation and Continuous Improvement, Wiley-IEEE Computer Society Press, 2008.
\item C. Seybold and R. Keller, Aligning Software Maintenance to the Offshore Reality, European Conference on Software Maintenance and Reengineering, 2008.
\item IEEE Std 2675-2021, IEEE Standard for DevOps: Building Reliable and Secure Systems Including Application Build, Package and Deployment, 2021.
\item W. Titus, T. Manshreck, and H. Wright, Software Engineering at Google: Lessons Learned from Programming over Time, O’Reilly Media, 2020.
\item A. Abran and H. Nguyenkim, Measurement of the Maintenance Process from a Demand-Based Perspective, Journal of Software Maintenance: Research and Practice, 1993.
\item M. M. Lehman, Laws of Software Evolution Revisited, European Workshop on Software Process Technology, 1996.
\item J. Humble and D. Farley, Continuous Delivery: Reliable Software Releases through Build, Test, and Deployment Automation, Addison-Wesley, 2010.
\item ISO/IEC/IEEE 12207:2017, Systems and software engineering — Software life cycle processes, 2017.
\end{itemize}
\section{9 画像 TODO 一覧}
本文に配置した画像 TODO を再掲する。実際に図を作る場合は、各 TODO の画像生成用プロンプトをそのまま使うと、A4 資料に合う図を作りやすい。\par
\begin{itemize}
\item 07 章の全体地図。
\item ソフトウェア進化の循環図。
\item ソフトウェア保守分類図。
\item 影響分析の連鎖図。
\item 技術的負債と保守費用の関係図。
\item ソフトウェア保守プロセス図。
\item 保守計画の四層図。
\item CI/CD/テスト/デプロイのパイプライン図。
\item 保守可視化ダッシュボード。
\item 保守ツールの分類図。
\end{itemize}
\section{10 章末確認問題}
\begin{itemize}
\item ソフトウェア保守は、なぜ単なるバグ修正ではないと言えるか。
\item 是正保守、予防保守、適応保守、追加保守、完全化保守の違いを説明せよ。
\item 緊急保守は、是正保守とどのように違うか。
\item 保守担当者の「限られた理解」は、なぜ保守費用を増やすのか。
\item 回帰テストが保守で重要になる理由を説明せよ。
\item 影響分析では、どのような情報を確認すべきか。
\item 保守性が低いソフトウェアは、どのように技術的負債を生むか。
\item 開発チームと保守チームを同一にする場合と分離する場合の利点と注意点を説明せよ。
\item 継続的統合、継続的デリバリ、継続的テスト、継続的デプロイは、保守にどう役立つか。
\item 保守ツールの中から三つを選び、何を支援するか説明せよ。
\end{itemize}
\boxblock{解答の方向性}{保守は、運用中のソフトウェアを支援し、変更し、価値を保つ活動である。分類問題では、欠陥修正か、潜在欠陥の予防か、環境適応か、機能追加か、品質改善かを見分ける。影響分析では、変更対象、依存関係、テスト、性能、安全性、セキュリティ、費用、リスクを確認する。}
\end{document}